
\medskip
\noindent\emph{Proof of \SGARefLink{sga3:xx:1.7}{1.7}.}
Let \(\mathfrak m\) be the radical of \(A\), and put
\[
  S_n=\Spec(A/\mathfrak m^{n+1}),\qquad n\geq 0.
\]
We first show, by induction on \(n\), that \(p_0\) can be extended,
for every \(n\), to a morphism of \(S_n\)-groups
\[
  p_n\colon\mathbb G_{a,S_n}\longrightarrow G_{S_n},
\]
normalized by \(T_{S_n}\) with multiplier \(\alpha_{S_n}\), in such
a way that
\[
  p_{n+1}\times_{S_{n+1}}S_n=p_n.
\]

Let \(H=T\mathbin{\cdot_\alpha}\mathbb G_{a,S}\). Denote by \(f_n\)
the morphism \(H_{S_n}\to G_{S_n}\) defined by \(p_n\). We
provisionally admit the following lemma.

\paragraph{\textbf{Lemma 1.9.}}
\SGARefTarget{sga3:xx:1.9}
\emph{Let \((G,T,\alpha)\) be a \(k\)-elementary system, where \(k\)
is algebraically closed, and let
\[
  p\colon\mathbb G_{a,k}\longrightarrow G
\]
be a monomorphism normalized by \(T\), with multiplier \(\alpha\).
Then}
\[
  \mathrm H^2\bigl(T\mathbin{\cdot_\alpha}\mathbb G_{a,k},
    \mathfrak g\bigr)=0.
\]
\emph{Here \(T\mathbin{\cdot_\alpha}\mathbb G_{a,k}\) acts on
\(\mathfrak g\) through the morphism
\(T\mathbin{\cdot_\alpha}\mathbb G_{a,k}\to G\) defined by \(p\)
and through the adjoint representation of \(G\).}

By \SGARefLink{sga3:iii:corollary:2-8}{Exposé III, 2.8}, \(f_n\) therefore extends to a morphism of
\(S_{n+1}\)-groups
\[
  f'_{n+1}\colon H_{S_{n+1}}\longrightarrow G_{S_{n+1}}.
\]
The morphism \(f'_{n+1}\) and the canonical immersion of
\(T_{S_{n+1}}\) into \(G_{S_{n+1}}\) have the same restriction to
\(T_{S_n}\). By \SGARefLink{sga3:iii:corollary:2-5}{Exposé III, 2.5}, there is an element
\(g\in G(S_{n+1})\) such that
\[
  g\times_{S_{n+1}}S_n=e
\]
and such that
\[
  f_{n+1}=\intaut(g)\circ f'_{n+1}
\]
restricts to \(T_{S_{n+1}}\) by the canonical immersion. Let
\(p_{n+1}\) be the restriction of \(f_{n+1}\) to
\(\mathbb G_{a,S_{n+1}}\). It is normalized by \(T_{S_{n+1}}\),
has multiplier \(\alpha_{S_{n+1}}\), and extends \(p_n\).

We have thus constructed a coherent system \((f_n)\); it remains to
algebraize it. For this we use:

\paragraph{\textbf{Lemma 1.10.}}
\SGARefTarget{sga3:xx:1.10}
\emph{Let \(A\) be a complete noetherian local ring,
\(\mathfrak m\) its maximal ideal,
\[
  S=\Spec(A),\qquad S_n=\Spec(A/\mathfrak m^{n+1}),
\]
let \(T\) be an \(S\)-torus, let \(\alpha\) be a nonzero character
of \(T\), and let \(X\) be an affine \(S\)-scheme on which \(T\)
acts. Let \(T\) act on \(\mathbb G_{a,S}\) through \(\alpha\).
Let \(q\geq0\), and let \((f_n)_{n\geq0}\) be a coherent system of
morphisms}
\[
  f_n\colon\mathbb G_{a,S_n}^{\,q}\longrightarrow X_{S_n}
\]
\emph{of objects with operators \(T_{S_n}\). There exists a unique
morphism of objects with operators \(T\)}
\[
  f\colon\mathbb G_{a,S}^{\,q}\longrightarrow X
\]
\emph{that induces the \(f_n\) \textup{(compare \SGARefLink{sga3:ix:theorem:7:1}{Exposé IX, 7.1})}.}

\paragraph{\textbf{Corollary 1.11.}}
\SGARefTarget{sga3:xx:1.11}
\emph{If \(X\) is a group with operator group \(T\), and if the
\(f_n\) are group morphisms, then \(f\) is also a group morphism.}

Indeed, apply the uniqueness assertion of the lemma to the two
morphisms
\(\mathbb G_{a,S}^{\,2q}\to X\) deduced from \(f\) in the usual
way.

\medskip
\noindent\emph{Proof of \SGARefLink{sga3:xx:1.10}{1.10}.}
Suppose that \(T\) is split, as it is in the application of \SGARefLink{sga3:xx:1.10}{1.10} to
the proof of \SGARefLink{sga3:xx:1.5}{1.5}. By \SGARefLink{sga3:i:corollary:4-7-3-1}{Exposé I, 4.7.3.1}, the functor
\[
  X\longmapsto\mathscr A(X)
\]
defines an equivalence between the category of affine \(S\)-schemes
equipped with an action of \(T\) and the category opposite to that
of \(M\)-graded \(S\)-algebras, where
\[
  M=\operatorname{Hom}_{S\text{-gr.}}(T,\mathbb G_{m,S}).
\]
Thus there are gradings
\[
  B=\mathscr A(X)=\bigoplus_{m\in M}B_m,
  \qquad
  C=\mathscr A(\mathbb G_{a,S}^{\,q})
    =\bigoplus_{m\in M}C_m.
\]
Each \(C_m\) is visibly a finite free \(A\)-module. Indeed,
\(C_m=0\) if \(m\) is not a multiple of \(\alpha\), whereas, if
\(m=d\alpha\), then \(C_m\) is isomorphic to the \(A\)-module of
homogeneous polynomials of degree \(d\) in \(q\) variables. Put
\[
\begin{aligned}
  \widehat B_m
    &=\varprojlim_n B_m\otimes_A(A/\mathfrak m^{n+1}),\\
  \widehat C_m
    &=\varprojlim_n C_m\otimes_A(A/\mathfrak m^{n+1}),\\
  \widehat B&=\bigoplus_{m\in M}\widehat B_m,
  &\widehat C&=\bigoplus_{m\in M}\widehat C_m.
\end{aligned}
\]
There are canonical morphisms of \(M\)-graded algebras
\[
  g_B\colon B\longrightarrow\widehat B,
  \qquad
  g_C\colon C\longrightarrow\widehat C.
\]
The preceding observation shows that \(g_C\) is an isomorphism.
Giving a coherent system \((f_n)\) as in the statement is
equivalent to giving a morphism of graded \(A\)-algebras
\[
  \widehat F\colon\widehat B\longrightarrow\widehat C.
\]
Finding a morphism \(f\) as in the statement is equivalent to
finding a morphism of graded \(A\)-algebras \(F\colon B\to C\)
making the diagram
\[
\begin{tikzcd}[column sep=large,row sep=large]
B \arrow[r,"F"] \arrow[d,"g_B"']
  & C \arrow[d,"g_C"']\\
\widehat B \arrow[r,"\widehat F"]
  & \widehat C
\end{tikzcd}
\]
commute. Since \(g_C\) is an isomorphism, the existence and
uniqueness of \(F\) are immediate. This proves \SGARefLink{sga3:xx:1.10}{1.10}.

To complete the proof of \SGARefLink{sga3:xx:1.5}{1.5}, it remains only to prove \SGARefLink{sga3:xx:1.9}{1.9}.
